\chapter{Conclusão}\label{cap:conclusao}

Retomando o objetivo geral de avaliar experimentalmente a capacidade dos LLMs na resolução de problemas do Nível Sênior da OBI (Modalidade Programação), este estudo trouxe cinco cenários nos quais, dentro da configuração utilizada de temperatura zero ($T = 0{,}0$) e apenas uma requisição ($\passat{1}$), os modelos divergem quanto à abordagem de melhor desempenho. No cenário da PP1, o \textbf{DeepSeek v4.1 Flash} e o \textbf{Gemini 3.8 Flash} mostraram uma amplitude menor de variação para questões e \textit{prompts} puramente textuais. Em sequência, para elementos visuais (PP2), no \textbf{DeepSeek v4.1 Flash} houve degradação comparativamente aos outros dois modelos na linguagem de melhor desempenho geral (\linguagem{C++}); isso mostra que modelos multimodais nem sempre conseguem performar bem com imagens no formato Base64 (\modalidade{img}), observando-se que o desempenho melhora quando existem descrições textuais das imagens (\modalidade{description}). Na avaliação de restrições (PP3), verificou-se que na linguagem mais sensível (\linguagem{C}) os modelos acabaram tratando como ruído informações que são importantes para os competidores da OBI (Modalidade Programação, Nível Sênior), como o limite de memória e o limite de tempo. Nesse caso de linguagens que precisam de maior cuidado na geração de código (PP4), notou-se que a aplicação de técnicas de \textit{prompt} é uma boa alternativa, destacando-se o \textit{Few-Shot} como melhor custo-benefício. Além disso, na comparação geracional (PP5), mostrou-se que os modelos de uma mesma família estão melhorando o custo-benefício e a geração de código para cenários de programação competitiva.

Por consequência, este Trabalho de Conclusão de Curso trouxe contribuições como uma base estruturada e validada por oráculo de 125 problemas do Nível Sênior da OBI (2016--2025) com 4.681 casos de teste, uma proposta de arquitetura experimental, uma ferramenta executável localmente e um \textit{benchmark} na resolução de problemas algorítmicos, alinhado às referências da literatura~\citep{chen2021evaluating,li2022competition,zheng2025livecodebench,alcantara2025assessing}. Para o cenário educacional, utilizar LLMs como motor de raciocínio em sistemas tutores inteligentes exige verificar se os modelos conseguem auxiliar os alunos, e não substituí-los, como mostram os estudos de~\citet{priyanka2025development,peng2025enhancing}.

Assim, como trabalhos futuros, a ferramenta pode ser expandida para avaliações em outros tipos de \textit{datasets} de problemas de programação competitiva, com perspectivas de utilização na geração de código em relação à métrica $\passat{k}$ e variando a temperatura, buscando avaliar a ``criatividade'' dos modelos. Além disso, vislumbra-se incorporar mecanismos de autocorreção para códigos já produzidos anteriormente, podendo evoluir para um sistema multiagente que utilize a metodologia de Pólya --- com cada agente dedicado a uma etapa do ciclo de resolução de problemas (Compreensão, Planejamento, Codificação (Execução) e Teste) ---, comunicando-se entre si por meio de protocolos bem definidos.
